%
%  $Description: Author guidelines and sample document in LaTeX 2.09$
%
%  $Author: ienne $
%  $Date: 1995/09/15 15:20:59 $
%  $Revision: 1.4 $
%

\documentclass[times, 10pt,twocolumn]{article}
\usepackage{latex8}
\usepackage{times}

%\documentstyle[times,art10,twocolumn,latex8]{article}

%-------------------------------------------------------------------------
% take the % away on next line to produce the final camera-ready version
\pagestyle{empty}

%-------------------------------------------------------------------------
\begin{document}

\title{Design and Implementation of Distributed Applications\\
dida-tkvs}

\author{Gonçalo Rua, João Gouveia, Francisco Salgueiro\\
  Instituto Superior Técnico\\ Av. Rovisco Pais 1\\ \{goncalonrua,joaotalonegouveia,francisco.salgueiro\}@tecnico.ulisboa.pt\\
}

\maketitle
\thispagestyle{empty}

\begin{abstract}
    dida-tkvs is a replicated key-value store composed of multiple servers,
    transactional clients, and a configuration console. High-performance
    consistency across replicas is achieved using the Multi-Paxos consensus
    algorithm. In this paper, we focus on the most critical aspects of our implementation.
\end{abstract}


%-------------------------------------------------------------------------
\Section{Introduction}
    In distributed systems, ensuring consistency and coordination across
    multiple nodes is a challenging task, especially in environments where
    failures can occur. Paxos is a consensus algorithm that addresses
    these issues by allowing a set of nodes to agree on a single value, even
    in the presence of failures. In this project, we implemented a transactional
    key-value store replicated using Paxos to ensure that all servers maintain a
    consistent view of the data.

    The system we developed operates in phases, starting from basic client-server
    communication and evolving into a fully reconfigurable Paxos-based system.
    This paper highlights the key components of the system, focusing on the main loop
    that drives the consensus process and handles new client requests. The logic behind
    selecting the next round number for the leader is also discussed, showcasing how
    our system ensures liveness and consistency.


%-------------------------------------------------------------------------
\Section{System Overview}
    The key-value store consists of multiple servers, each holding a replicated state
    of the store. Clients initiate transactions, which involve reading and updating
    key-value pairs. Servers coordinate through Paxos to assign unique timestamps to
    updates, ensuring that conflicting operations are managed correctly.
    The system can dynamically reconfigure itself, allowing different sets of servers
    to take on the roles of proposers and acceptors while ensuring that all nodes act
    as learners.

%-------------------------------------------------------------------------
\SubSection{Main Loop for Transaction Proposals and Execution}
    The main loop is central to the system’s operation, where the leader proposes
    transactions, and the learners execute them. The flow is as follows:

\begin{itemize}
    \item Transaction Proposal by Leader: To be completed.
    \item Transaction Execution by Learners: To be completed.
\end{itemize}

%-------------------------------------------------------------------------
\SubSection{Handling New Requests}
    When a new request is received from a client, the system must ensure that the
    transaction is processed in a consistent and timely manner. This involves
    coordinating between the leader and the learners to ensure the request is either
    committed or aborted based on the current state of the key-value store. To be completed.

%-------------------------------------------------------------------------
\SubSection{Computing the Next Round Number (rnd) for the Leader}
    The leader plays a crucial role in managing the order of transactions.
    One of its key responsibilities is determining the next round number (rnd)
    to use in Paxos. This ensures that no conflicting proposals are made,
    and that progress is always possible. To be completed.

%-------------------------------------------------------------------------
\Section{Conclusion}
    The implementation of our transactional key-value store using Paxos
    demonstrates the effectiveness of the algorithm in maintaining consistency
    across distributed servers. The main loop and its handling of transaction
    proposals and executions are vital to the system's overall functionality.
    While we have discussed key aspects of the system, including how the leader
    computes the next round number, further optimization could be explored,
    especially in handling client requests efficiently. Overall, our implementation
    provides a robust foundation for distributed consensus in transactional systems.

\begin{figure}[h]
   \caption{Example of caption.}
\end{figure}

%-------------------------------------------------------------------------
\nocite{ex1,ex2}
\bibliographystyle{latex8}
\bibliography{latex8}

\end{document}
